\documentclass[11pt, a4paper]{article}

% --- UNIVERSAL PREAMBLE BLOCK ---
\usepackage[a4paper, top=2.5cm, bottom=2.5cm, left=2cm, right=2cm]{geometry}
\usepackage{fontspec}
\usepackage[english, bidi=basic, provide=*]{babel}

% Set default font to Noto Sans
\babelfont{rm}{Noto Sans}

\usepackage{amsmath}
\usepackage{amssymb}
\usepackage{booktabs}
\usepackage{xcolor}

\title{Foundations of Prompt Arithmetic}
\author{Formal Logic Systems}
\date{\today}

\begin{document}

\maketitle

\begin{abstract}
This document introduces the formal notation and operational logic for Prompt Arithmetic. By treating natural language instructions as algebraic variables, we can define operations such as addition, subtraction, multiplication, division, exponentiation, and vector scoping to predict and control the output behavior of Large Language Models (LLMs).
\end{abstract}

\section{Introduction}
Prompt arithmetic is the study of how combining discrete sets of instructions (prompts) influences the resulting output space. In this framework, a prompt $p$ is defined as a functional operator that maps an input state to a specific linguistic or programmatic response. 

The practical execution of this algebra is managed by semantic compiler engines (such as \textbf{Glupe}), which act as the Prompt Virtual Machine (PVM), while build agents (such as \textbf{Juan}) orchestrate the dependency graphs, vector scoping, and file I/O.

\section{Prompt Addition}
The sum of two prompts is defined as the concatenation of their semantic intents.
$$r = p + q$$
Example: If $p \rightarrow \text{print ``hi''}$ and $q \rightarrow \text{print ``there''}$, then $p+q \rightarrow \text{print ``hi there''}$.

\subsection{Instruction Conflicts and Compile-Time Errors}
A critical challenge in prompt addition is the occurrence of \textbf{semantic contradictions}.
Consider the following prompts:
\begin{align*}
    \$: x &\rightarrow \text{print ``hi''} \\
    \$: y &\rightarrow \text{do not print anything}
\end{align*}
When we define $\$z = \$x + \$y$, we encounter a logical impasse. In a rigorous DSL, this results in a \textbf{Compile-Time Error} $\bot$.
$$ \text{If } \text{Intent}(x) \cap \text{Intent}(y) = \emptyset, \text{ then } x + y = \bot $$

\section{Prompt Subtraction}
Prompt subtraction is the operation of removing specific instruction members or constraints from a base prompt, acting as an \textbf{Override} mechanism.

\subsection{Case 1: Simple Subtraction (Subset Removal)}
When a prompt $y$ is a subset of the logic contained in $x$, subtraction removes those shared members.
\begin{align*}
    \$: x &\rightarrow \text{call API and print ``hi''} \\
    \$: y &\rightarrow \text{call API} \\
    \$: z &= \$x - \$y \rightarrow \text{print ``hi''}
\end{align*}
Here, the instruction to ``call API'' is subtracted, leaving only the residual instruction.

\subsection{Case 2: Subtraction with Remainder (Negative Prompts)}
When the subtrahend contains logic not present in the minuend, we encounter a ``negative remainder.'' A negative prompt state $-(\text{logic})$ serves as an \textbf{Inhibitor}.
\begin{align*}
    \$: a &= \$y - \$x \\
    \$: a &= (\text{call API}) - (\text{call API and print ``hi''}) \\
    \$: a &= -(\text{print ``hi''})
\end{align*}
While positive logic dictates what the model \textit{must do}, negative logic dictates what the model \textit{must suppress}.

\subsection{The Null Condition}
The validity of the negative remainder is proven by the satisfaction of the \textbf{Null Condition}. When a positive instruction and its negative counterpart are added, they result in \texttt{NULL}.
\begin{align*}
    \$a + \$z &= (\$y - \$x) + (\$x - \$y) = \text{NULL} \\
    -(\text{print ``hi''}) &+ (\text{print ``hi''}) = \text{NULL}
\end{align*}

\section{Prompt Multiplication}
Prompt multiplication ($p \times q$) represents the \textbf{Iterative Execution} or \textbf{Mapping} of one prompt over the domain of another (Horizontal Scaling).
$$ \text{If } p \rightarrow \text{Action}, q \rightarrow \text{Iterator}, \text{ then } p \times q \rightarrow \text{Loop}(q, p) $$

\subsection{Example: Operational Loop}
Consider the following definitions:
\begin{align*}
    \$: p &\rightarrow \text{call API} \\
    \$: q &\rightarrow \text{pop stack} \\
    \$: r &= \$p \times \$q
\end{align*}
The multiplication yields a loop logic: ``Call API for every item popped from the stack.'' The prompt $q$ acts as the \textbf{Iterator}, while $p$ acts as the \textbf{Function Body}.

\subsection{Case: Scalar Multiplication}
If $q$ is a scalar $k$ (a raw integer), then $k \times p$ simply repeats the instruction $k$ times to increase its weight or recurrence in the LLM context window.

\section{Prompt Inversion and Division}
The introduction of the multiplicative inverse ($p^{-1}$ or $1/p$) allows for the formal definition of \textbf{Semantic Reversal} and \textbf{Conditional Disconnects}.

\subsection{The Multiplicative Inverse}
The inverse of a prompt $x$ is defined as the negation or undoing of the state intended by $x$.
\begin{align*}
    \$: x &\rightarrow \text{send email to client X} \\
    \$: y &= 1/\$x \rightarrow \text{do not send email to client X} \\
    \$: z &= 1/\$y = 1/(1/\$x) \rightarrow \text{send email to client X}
\end{align*}
Algebraically, $x \times (1/x) = \text{Identity}$, where Identity represents the state of not performing the specific action $x$ while maintaining the potential for other operations.

\subsection{Prompt Division as Conditional Undo}
Division ($y/x$) or the product of an action and an inverse ($x^{-1} \cdot y$) creates a logic where the reversal of $x$ is bound to the execution of $y$.
\begin{align*}
    \$: x &\rightarrow \text{connect to wifi} \\
    \$: y &\rightarrow \text{call API} \\
    \$: z &= \frac{1}{\$x} \cdot \$y = \text{disconnect from wifi for each API call}
\end{align*}
This defines a \textbf{Resource Cleanup Pattern}. The inverse prompt $1/x$ (disconnect) is mapped across the iterator/action $y$ (API call), ensuring that for every instance of the second, the reversal of the first is enforced.

\section{Fundamental States: The Zero and The Identity}
To complete the algebraic foundation, we must define the constants of the prompt vector space: the ``Zero'' and the ``One.'' These govern the baseline behaviors of the AI model.

\subsection{The Null State ($\emptyset$ or $0$)}
The Null State represents the complete absence of instruction or constraint. It is the \textbf{Additive Identity}.
\begin{align*}
    \$: \emptyset &\rightarrow \text{[Empty Instruction / Await Input]} \\
    \$p + \emptyset &= \$p
\end{align*}
In autonomous systems, if an agent's active prompt equation resolves to $\emptyset$ (e.g., through successive prompt subtractions), it acts as a \textbf{Halt Command}. It returns the LLM to its base, unsteered pre-training distribution, meaning there are no active constraints remaining in the context window. Furthermore, any prompt multiplied by zero yields the null state ($p \times 0 = \emptyset$).

\subsection{The Identity Element ($\mathbb{I}$ or $1$)}
The Identity Element is the \textbf{Multiplicative Identity}. It represents the state of pure preservation or verbatim execution.
\begin{align*}
    \$: \mathbb{I} &\rightarrow \text{Execute unaltered / Output input verbatim} \\
    \$p \times \mathbb{I} &= \$p
\end{align*}
This is crucial for evaluating model alignment: if a model applies an unintended style (such as unsolicited safety lectures or excessive politeness) to a neutral prompt, it is failing to maintain the Identity State.

\section{Properties of the Operators}
\begin{enumerate}
    \item \textbf{Additive Identity (Zero):} $p + \emptyset = p$.
    \item \textbf{Multiplicative Identity (One):} $p \times \mathbb{I} = p$.
    \item \textbf{Additive Inverse:} $p + (-p) = \emptyset$ (Total suppression).
    \item \textbf{Multiplicative Inverse:} $p \times (1/p) = \mathbb{I}$ (Action Neutralization).
    \item \textbf{Distributivity:} $p \times (q + s) = (p \times q) + (p \times s)$.
\end{enumerate}

\section{Prompt Exponentiation (Recursive Composition)}
Building upon prompt multiplication (defined as an iterative map), we can define exponentiation ($p^n$) as the \textbf{Recursive Composition} of a prompt upon its own output space (Vertical Refinement).

\subsection{Definition: The Self-Referential Loop}
If $p \times q$ dictates that the action $p$ is mapped over the domain of $q$, then $p \times p$ (or $p^2$) dictates that action $p$ is mapped over the output generated by the first instance of $p$.
$$ \$p^n = \underbrace{\$p \times \$p \times \dots \times \$p}_{n \text{ times}} $$
This operation mathematically models \textbf{Self-Correction}, \textbf{Iterative Refinement}, and \textbf{Meta-Cognition} within an autonomous LLM framework.

\subsection{Example: Iterative Refinement}
Consider an instruction designed for evaluating and rewriting a block of text:
\begin{align*}
    \$: p &\rightarrow \text{critique this draft and apply the feedback to rewrite it} \\
    \$: p^2 &= \$p \times \$p \rightarrow \text{rewrite the draft, then critique and rewrite the new version} \\
    \$: p^n &\rightarrow \text{perform } n \text{ cycles of self-critique and revision}
\end{align*}
By raising a prompt to a power $n$, the system instructs an agent to perform an $n$-depth recursive loop. This progressively distills or refines the output state, moving the response closer to a mathematical asymptote of quality or compression, shifting the developer from micromanaging specific improvements to commanding the depth of cognitive cycles.

\section{Prompt Equations and Scalar Addition}
By treating prompts as algebraic variables, we can construct formal equations to either synthesize new prompts or solve for missing logical constraints.

\subsection{Scalar Addition (Detail Increments)}
In Prompt Arithmetic, we define the addition of a scalar $k$ (where $k \in \mathbb{Z}$) to a prompt $p$ as a \textbf{Detail Increment}. It dictates that the LLM must execute $p$ and append $k$ standard units of elaboration (such as an example, a proof step, or a concrete application).

Consider the equation:
\begin{align*}
    \$: z &\rightarrow \text{solve PDE efficiently} \\
    \$x &= \$z + 1
\end{align*}
Here, $\$x$ evaluates to the synthesized state: \textit{``solve PDE efficiently, and provide 1 standard elaboration (e.g., one worked example).''}

\subsection{Solving for Unknown Prompts}
The true utility of prompt equations lies in determining the necessary prompt to transition an LLM from an initial state to a desired target state. 
$$\$A + \$x = \$B \implies \$x = \$B - \$A$$

\textbf{Example:}
\begin{align*}
    \$: A &\rightarrow \text{write a sorting algorithm} \\
    \$: B &\rightarrow \text{write a sorting algorithm with O(n log n) complexity} \\
    \$: x &= \$B - \$A \rightarrow \text{ensure O(n log n) complexity}
\end{align*}
By algebraically solving for $\$x$, an automated system can dynamically generate the exact prompt modifiers required to steer an LLM's output without rewriting the base prompt.

\section{High-Order Applications of Prompt Algebra}
The formalization of Prompt Arithmetic enables complex architectural patterns that were previously non-deterministic or computationally prohibitive.

\subsection{Dynamic Logic Extraction (Prompt Diffing)}
Developers can isolate the exact ``delta'' between two complex behaviors to create a portable logic vector:
$$ \$Vibe_{pirate} = \$Response_{pirate} - \$Response_{standard} $$
This extracted vector can then be added to an unrelated technical prompt to transfer stylistic or logical constraints with mathematical precision.

\subsection{Context Garbage Collection ($O(1)$ Agent State)}
Prompt Algebra treats context as a dynamic state equation. When a task is completed, an autonomous agent executes:
$$ S_{t+1} = S_t - \$Task_{completed} + \$Task_{next} $$
By algebraically removing completed constraints, the agent's context window remains $O(1)$, avoiding linear $O(n)$ degradation.

\subsection{Cross-Model Translation Constants}
Different models respond differently to the same text. If we know the baseline evaluation of a prompt $\$p$ on Model A, and we want to achieve the exact same semantic output space on Model B, we can solve for a Translation Constant ($k$):
$$ k_B \times \$p = \text{Output}_A $$

\section{Lexical Scoping and Prompt Vectors}
As Prompt Algebra scales to enterprise applications, the risk of "Context Poisoning" becomes the primary limiting factor. To solve this, we introduce \textbf{Prompt Vectors} and \textbf{Semantic Capture Lists}.

\subsection{The Semantic Closure}
Borrowing from the lambda calculus and modern programming languages, a prompt vector explicitly defines its dependencies before execution. 
Consider the syntax:
$$ \$: x \rightarrow y, z\ \{ \text{call api, clean buffer, } y^2, z^2 \} $$

In this notation, the vector $x$ captures \textit{only} the variables $y$ and $z$ from the global scope. The semantic engine (Glupe) isolates the execution environment inside the $\{ \dots \}$ block. The LLM is mathematically blocked from accessing any state outside of $y$ and $z$. This strict scoping ensures $O(1)$ garbage collection for the LLM's localized context.

\subsection{Prompt Matrices and SIMD Execution}
By defining prompts as isolated vectors, we unlock multi-dimensional synthesis. If we define a vector of actions and a vector of data domains:
\begin{align*}
    \$: \vec{A} &\rightarrow \{ \text{Create}, \text{Read}, \text{Delete} \} \\
    \$: \vec{D} &\rightarrow \{ \text{User}, \text{Order} \}
\end{align*}
We can calculate the \textbf{Dot Product} to synthesize an entire API matrix:
$$ \$: Matrix = \vec{A} \cdot \vec{D} $$

Because the scope of each operation is strictly encapsulated, a build agent (like Juan) can execute this matrix using \textbf{SIMD (Single Instruction, Multiple Data)} architecture, spinning up independent parallel threads without context bleeding.

\section{Hardware for Thought: Implementation via Semantic Containers}
Theoretical algebras require a physical or virtual machine to execute them. This execution environment is realized through \textbf{Semantic Containers} compiled by the Glupe engine. 

These containers act as isolated compiler blocks where prompt equations are solved. The Glupe engine verifies the algebraic integrity (catching $\bot$ contradictions), while the Juan build system maps the resolved deterministic output into the host environment's file structure. This "hybrid determinism" enables the synthesis of massive, production-grade applications from a single source file of intent.

\section{Formal Proof of State Contradiction}
To demonstrate the necessity of the compile-time error ($\bot$), we provide a formal proof of contradiction when adding mutually exclusive prompt vectors. 

\textbf{Theorem:} If prompt $p$ and prompt $q$ demand mutually exclusive output states, their sum $p + q$ must resolve to $\bot$ before inference begins.

\textbf{Proof:} \\
Let $\mathcal{S}$ be the universal set of all possible generated outputs from a language model. \\
Let an evaluation function $E(p)$ return the subset of $\mathcal{S}$ that completely satisfies the constraints of a given prompt $p$. 

By definition, the prompt addition $\$r = \$p + \$q$ requires the final output to satisfy the constraints of both $p$ and $q$ simultaneously. Therefore, the valid solution space for $r$ is the intersection of their individual spaces:
$$E(r) = E(p) \cap E(q)$$

Consider the following mutually exclusive prompt definitions:
\begin{align*}
    \$: p &\rightarrow \text{``output exactly 10 words''} \\
    \$: q &\rightarrow \text{``output exactly 50 words''}
\end{align*}

The valid solution spaces for these prompts are:
\begin{align*}
    E(p) &= \{ s \in \mathcal{S} \mid \text{word\_count}(s) = 10 \} \\
    E(q) &= \{ s \in \mathcal{S} \mid \text{word\_count}(s) = 50 \}
\end{align*}

Evaluating the intersection for the combined prompt $r$:
$$E(p) \cap E(q) = \emptyset$$

Because the solution space is the empty set ($\emptyset$), there exists no valid output string $s$ that can satisfy the equation. If an inference engine attempts to evaluate $r$, it will be forced into a probabilistic guess, necessarily violating at least one constraint and leading to hallucination.

To guarantee deterministic behavior and system integrity, the prompt compiler must intercept any operation yielding an empty intersection and substitute the $\bot$ (Halt/Error) operator:
$$\text{Since } E(p) \cap E(q) = \emptyset, \text{ then } \$p + \$q = \bot \quad \blacksquare$$

\section{Conclusion}
The formalization of Prompt Arithmetic represents a critical paradigm shift in artificial intelligence orchestration, transitioning the field of prompt engineering from a stochastic, heuristic-driven craft into a rigorous, deterministic branch of computer science. By establishing a comprehensive algebraic framework—encompassing addition, subtraction, iterative multiplication, conditional division, scalar equations, and explicit vector operations—we have demonstrated that natural language instructions can be treated as computable, invertible structures within a defined semantic space.

This mathematical approach directly addresses some of the most persistent challenges in modern AI deployment, particularly the inherent unpredictability of large language models. Through the introduction of the Identity Element ($\mathbb{I}$), the Null State ($\emptyset$), and formal proofs of state contradiction ($\bot$), we establish the theoretical groundwork for \textit{Prompt Compilers}: intermediary systems capable of optimizing, validating, and mathematically resolving instruction conflicts before inference even begins. 

Furthermore, as the industry advances toward highly autonomous agentic workflows and complex multi-model pipelines, the necessity for robust state management becomes paramount. Prompt Algebra provides the exact mechanics required to maintain long-term instruction fidelity. Executed through the Glupe semantic engine and orchestrated by the Juan build system, it allows autonomous architectures to dynamically add, subtract, and scale behavioral constraints using parallel SIMD workflows, without suffering from context window degradation, logic loops, or semantic drift. 

Ultimately, the principles outlined in this document move beyond theoretical curiosity; they lay the mathematical foundation necessary for building the next generation of safe, verifiable, and highly predictable artificial intelligence systems. By holding models to algebraic standards, we ensure that the future of human-AI interaction is governed by logic rather than chance.

\end{document}